% GENERATED FILE. Edit canonical lessons, then run npm run generate:books.
% canonical-source-hash: fnv1a64:fe37ed00d3b8b9c1
% canonical-lessons: AR-C63-rama, AR-C63-tara, AR-C63-sabaha, AR-C63-sanaa, AR-C63-dakhala

\chapter{To throw, To fly, To swim, To make, To enter}
\label{ch:ar-to-throw-to-fly-to-swim-to-make-to-enter}
\hlchaptermodality{63}

\begin{chapteropening}
\textbf{By the end of this chapter:} \emph{I can say to throw, to fly, to swim, to make and to enter.}

Complete the last lesson of chapter 63: I can say to throw, to fly, to swim, to make and to enter.
\end{chapteropening}

\section[\texorpdfstring{\ar{رمى} (ramā)}{رمى (ramā)}]{\ar{رمى} (ramā) — to throw (he threw)}

\label{lesson:AR-C63-rama}



\begin{quote}
Before the new one: say the Arabic for to wear, then the Arabic for to leave.
\end{quote}

\subsection*{You'll want to know: \ar{رمى}}
\textbf{\ar{رمى}} — \emph{ramā} — \textquotedblleft{}to throw (he threw)\textquotedblright{}.

\begin{cousinweb}
\textbf{\ar{رمى}} (\emph{ramā}) is built on the native Arabic root \emph{r-m-y}, \textquotedblleft{}to throw, to shoot\textquotedblright{} (Lane). \emph{Ramā} is the plain verb, \textquotedblleft{}he threw\textquotedblright{}.
\end{cousinweb}

\begin{grammarlens}[title={What you've built}]
One more word for this chapter.
\end{grammarlens}

\subsection*{Your turn}
\begin{itemize}
\raggedright
  \item \emph{Say it:} \emph{ramā}
  \item \emph{Say it:} \emph{ramā}, once more
  \item \emph{Say it:} say \emph{ramā}
  \item \emph{Recall:} say the Arabic for to wear, then the Arabic for to leave, then say \emph{ramā} again
\end{itemize}

\begin{tcolorbox}[breakable,colback=teal!4,colframe=teal!35!black,title={Before you move on}]
What does \ar{رمى} mean? (\textquotedblleft{}To throw (he threw)\textquotedblright{}.) Say it once more.
\end{tcolorbox}
\section[\texorpdfstring{\ar{طار} (ṭāra)}{طار (ṭāra)}]{\ar{طار} (ṭāra) — to fly (he flew)}

\label{lesson:AR-C63-tara}



\begin{quote}
Before the new one: say the Arabic for to leave, then the Arabic for to throw.
\end{quote}

\subsection*{You'll want to know: \ar{طار}}
\textbf{\ar{طار}} — \emph{ṭāra} — \textquotedblleft{}to fly (he flew)\textquotedblright{}.

\begin{cousinweb}
\textbf{\ar{طار}} (\emph{ṭāra}) is built on the native Arabic root \emph{ṭ-y-r}, \textquotedblleft{}to fly\textquotedblright{} (Lane). The same root gives \emph{ṭāʾir}, \textquotedblleft{}a bird\textquotedblright{}, and \emph{maṭār}, \textquotedblleft{}an airport\textquotedblright{}.
\end{cousinweb}

\begin{grammarlens}[title={What you've built}]
One more word for this chapter.
\end{grammarlens}

\subsection*{Your turn}
\begin{itemize}
\raggedright
  \item \emph{Say it:} \emph{ṭāra}
  \item \emph{Say it:} \emph{ṭāra}, once more
  \item \emph{Say it:} say \emph{ṭāra}
  \item \emph{Recall:} say the Arabic for to leave, then the Arabic for to throw, then say \emph{ṭāra} again
\end{itemize}

\begin{tcolorbox}[breakable,colback=teal!4,colframe=teal!35!black,title={Before you move on}]
What does \ar{طار} mean? (\textquotedblleft{}To fly (he flew)\textquotedblright{}.) Say it once more.
\end{tcolorbox}
\section[\texorpdfstring{\ar{سبح} (sabaḥa)}{سبح (sabaḥa)}]{\ar{سبح} (sabaḥa) — to swim (he swam)}

\label{lesson:AR-C63-sabaha}



\begin{quote}
Before the new one: say the Arabic for to throw, then the Arabic for to fly.
\end{quote}

\subsection*{You'll want to know: \ar{سبح}}
\textbf{\ar{سبح}} — \emph{sabaḥa} — \textquotedblleft{}to swim (he swam)\textquotedblright{}.

\begin{cousinweb}
\textbf{\ar{سبح}} (\emph{sabaḥa}) is built on the native Arabic root \emph{s-b-ḥ}, \textquotedblleft{}to swim; to glorify God\textquotedblright{} (Lane). Lane gives swimming and, in Form II, praising God under one root.
\end{cousinweb}

\begin{grammarlens}[title={What you've built}]
One more word for this chapter.
\end{grammarlens}

\subsection*{Your turn}
\begin{itemize}
\raggedright
  \item \emph{Say it:} \emph{sabaḥa}
  \item \emph{Say it:} \emph{sabaḥa}, once more
  \item \emph{Say it:} say \emph{sabaḥa}
  \item \emph{Recall:} say the Arabic for to throw, then the Arabic for to fly, then say \emph{sabaḥa} again
\end{itemize}

\begin{tcolorbox}[breakable,colback=teal!4,colframe=teal!35!black,title={Before you move on}]
What does \ar{سبح} mean? (\textquotedblleft{}To swim (he swam)\textquotedblright{}.) Say it once more.
\end{tcolorbox}
\section[\texorpdfstring{\ar{صنع} (ṣanaʿa)}{صنع (ṣanaʿa)}]{\ar{صنع} (ṣanaʿa) — to make (he made)}

\label{lesson:AR-C63-sanaa}



\begin{quote}
Before the new one: say the Arabic for to fly, then the Arabic for to swim.
\end{quote}

\subsection*{You'll want to know: \ar{صنع}}
\textbf{\ar{صنع}} — \emph{ṣanaʿa} — \textquotedblleft{}to make (he made)\textquotedblright{}.

\begin{cousinweb}
\textbf{\ar{صنع}} (\emph{ṣanaʿa}) is built on the native Arabic root \emph{ṣ-n-ʿ}, \textquotedblleft{}to make, to fashion\textquotedblright{} (Lane). The same root gives \emph{maṣnaʿ}, \textquotedblleft{}a factory\textquotedblright{}.
\end{cousinweb}

\begin{grammarlens}[title={What you've built}]
One more word for this chapter.
\end{grammarlens}

\subsection*{Your turn}
\begin{itemize}
\raggedright
  \item \emph{Say it:} \emph{ṣanaʿa}
  \item \emph{Say it:} \emph{ṣanaʿa}, once more
  \item \emph{Say it:} say \emph{ṣanaʿa}
  \item \emph{Recall:} say the Arabic for to fly, then the Arabic for to swim, then say \emph{ṣanaʿa} again
\end{itemize}

\begin{tcolorbox}[breakable,colback=teal!4,colframe=teal!35!black,title={Before you move on}]
What does \ar{صنع} mean? (\textquotedblleft{}To make (he made)\textquotedblright{}.) Say it once more.
\end{tcolorbox}
\section[\texorpdfstring{\ar{دخل} (dakhala)}{دخل (dakhala)}]{\ar{دخل} (dakhala) — to enter (he entered)}

\label{lesson:AR-C63-dakhala}



\begin{quote}
Before the new one: say the Arabic for to swim, then the Arabic for to make.
\end{quote}

\subsection*{You'll want to know: \ar{دخل}}
\textbf{\ar{دخل}} — \emph{dakhala} — \textquotedblleft{}to enter (he entered)\textquotedblright{}.

\begin{cousinweb}
\textbf{\ar{دخل}} (\emph{dakhala}) is built on the native Arabic root \emph{d-kh-l}, \textquotedblleft{}to enter\textquotedblright{} (Lane). \emph{Dakhala} is the plain verb, \textquotedblleft{}he entered\textquotedblright{}.
\end{cousinweb}

\begin{grammarlens}[title={What you've built}]
One more word for this chapter.
\end{grammarlens}

\subsection*{Your turn}
\begin{itemize}
\raggedright
  \item \emph{Say it:} \emph{dakhala}
  \item \emph{Say it:} \emph{dakhala}, once more
  \item \emph{Say it:} say \emph{dakhala}
  \item \emph{Recall:} say the Arabic for to swim, then the Arabic for to make, then say \emph{dakhala} again
\end{itemize}

\begin{tcolorbox}[breakable,colback=teal!4,colframe=teal!35!black,title={Before you move on}]
What does \ar{دخل} mean? (\textquotedblleft{}To enter (he entered)\textquotedblright{}.) Say it once more.
\end{tcolorbox}
